\documentclass[a4paper,10pt]{report}\input{../0.Préambule/Préambule_cours_prof}\begin{document}

\newpage
\setcounter{chapter}{4}
\chapter{Nombre en écriture fractionnaire}

\section{Vocabulaire}

\subsection{Fractions et quotients}

\begin{ex}
Le quotient de $3,5$ par $2$ est le résultat de la division de $3,5$ par $2$. \\On le note $3,5 \div 2 = \dfrac{3,5}{2} = 1,75$.

\noindent $\dfrac{3,5}{2}$ est l'écriture fractionnaire, $1,75$ est l'écriture décimale
\end{ex}

\vspace{-3.7cm}
\begin{flushright}
\begin{tikzpicture}[line cap=round,line join=round,>=triangle 45,x=1.0cm,y=1.0cm,scale=0.6]
\clip(0.,0.5) rectangle (14.,6.);
\draw [->,line width=1.2pt] (4.,5.) -- (6.,4.);
\draw [->,line width=1.2pt] (5.,2.) -- (7.,3.);
\draw [->,line width=1.2pt] (11.,5) -- (9.,4.2);
\draw [->,line width=1.2pt] (11.,2.) -- (9.,2.7);
\begin{scriptsize}
\draw[color=black] (7.2,3.4554545454545424) node {\Large{$a \div b = \dfrac{a}{b}$}};
\draw[color=black] (3.5,5.4554545454545424) node {\Large{Dividende}};
\draw[color=black] (4.26,1.4754545454545451) node {\Large{Diviseur}};
\draw[color=black] (11.02,5.635454545454542) node {\Large{Numérateur}};
\draw[color=black] (11.28,1.4154545454545453) node {\Large{Dénominateur}};
\end{scriptsize}
\end{tikzpicture}
\end{flushright}


\subsection{Fractions et produits}

\vspace{-3mm}
\noindent \begin{minipage}{8.5cm}
\begin{dft}
La fraction $\dfrac{\textcolor{red}{a}}{\textcolor{blue}{b}}$ est la solution de l'opération à trou : $$\textcolor{blue}{b} \times \dots = \textcolor{red}{a}$$
\end{dft}
\end{minipage}
\hspace{5mm}
\begin{minipage}{8cm}
\begin{ex}

%Le nombre qui multiplié par $7$ donne $8$ est $\dfrac{8}{7}$ . ( $\textcolor{blue}{7} \times \dfrac{\textcolor{red}{8}}{\textcolor{blue}{7}} = \textcolor{red}{8}$ )

\medskip
\noindent Le nombre $\dfrac{5}{6}$ est le nombre qui multiplié par $6$ donne $5$. ( $\textcolor{blue}{6} \times \dfrac{\textcolor{red}{5}}{\textcolor{blue}{6}} = \textcolor{red}{5}$ )
\end{ex}
\end{minipage}

\subsection{Nombre entier, nombre rationnel et nombre décimal}

\begin{ex}
\begin{enumerate}
\item[•] $\dfrac{35}{7} = 35 \div 7 = 5$\hspace{3.6cm}$\dfrac{35}{7}$ est un nombre rationnel qui est un nombre entier.
\item[•] $\dfrac{15}{4}= 15 \div 4 = 3,75$\hspace{3.2cm}$\dfrac{15}{4}$ est un nombre rationnel qui est un nombre décimal.
\item[•] $\dfrac{10}{3}= 10 \div 3 \approx 3,3333333 \dots$\hspace{2cm}$\dfrac{10}{3}$ n'est pas un nombre décimal car la division ne se termine pas.

\noindent On ne peut donc pas donner une valeur exacte du quotient $\dfrac{10}{3}$.

\noindent On ne peut en donner qu'une valeur approchée : $\dfrac{10}{3} \approx 3,33$ ( valeur approchée au centième par défaut)
\end{enumerate}
\end{ex}

\vspace{-3mm}
\section{Proportion}

\begin{ex}
Dans le mot FRACTION, $5$ lettres sur les $8$ sont des consonnes.
On dit que la proportion (ou la fréquence) de consonnes du mot FRACTION est $\dfrac{5}{8}$.

\medskip
\noindent $\dfrac{5}{8} = 5 \div 8 = 0,625 = \dfrac{62,5}{100}$ 

\noindent donc cette fréquence peut aussi s'exprimer par le nombre décimal $0,625$ ou par le pourcentage $62,5 \%$.
\end{ex}

\begin{dft}
\begin{enumerate}
\item[•] Lorsque le numérateur et le dénominateur sont des entiers, on parle de fraction.
\item[•] L'ensemble des nombres qui peuvent s'écrire $\dfrac{a}{b}$ où $a$ est un nombre relatif et $b$ un nombre relatif non nul est appelé l'ensemble des nombres rationnels.
\end{enumerate}
\end{dft}

\begin{ex}
\begin{enumerate}
\item[•] $\dfrac{4}{5}$ est une fraction. $4$ est le numérateur et $5$ est le dénominateur.
\item[•] $\dfrac{6,5}{8}$ n'est pas une fraction, c'est une écriture fractionnaire.
\end{enumerate}
\end{ex}

\begin{rmq}
Les nombres décimaux et les nombres relatifs sont des nombres rationnels. Par exemple, $-2,35 = \dfrac{-235}{100}$
\end{rmq}

\begin{regle}[Différents sens de l'écriture fractionnaire]
La fraction $\dfrac{7}{5}$ se lit \og sept cinquièmes \fg{}. Cette fraction est égale :
\begin{enumerate}
\item[•] à $7$ fois un cinquième car $\dfrac{7}{5} = 7 \times \dfrac{1}{5} = \dfrac{1}{5} \times 7$
\item[•] au quotient de $7$ par $5$ car $\dfrac{7}{5} = 7 \div 5$
\item[•] au nombre qui multiplié par $5$ donne $7$ car $7 = 5 \times \dfrac{7}{5} = \dfrac{7}{5} \times 5$
\item[•] au nombre $1 + \dfrac{2}{5}$
\end{enumerate}
\end{regle}

\medskip
On peut aussi représenter cette fraction de plusieurs façon, par exemple:

\noindent \begin{tikzpicture}[line cap=round,line join=round,>=triangle 45,x=1.0cm,y=1.0cm,scale=0.7]
\clip(3.,2.) rectangle (7.5,6.);
\draw [line width=1.6pt] (5.18,4.02) circle (1.8201098867925534cm);
\draw [line width=1.2pt] (7.,4.)-- (5.18,4.02);
\draw [line width=1.2pt] (5.18,4.02)-- (5.767189238011339,5.738460372547006);
\draw [line width=1.2pt] (3.7215252327692143,5.117446058978674)-- (5.18,4.02);
\draw [line width=1.2pt] (5.18,4.02)-- (3.68,2.98);
\draw [line width=1.2pt] (5.18,4.02)-- (5.7,2.28);
\draw [shift={(5.18,4.02)},line width=1.2pt,color=black,fill=blue,fill opacity=0.2]  (0,0) --  plot[domain=1.241539285032363:2.4965180680191064,variable=\t]({1.*1.8160113582383595*cos(\t r)+0.*1.8160113582383595*sin(\t r)},{0.*1.8160113582383595*cos(\t r)+1.*1.8160113582383595*sin(\t r)}) -- cycle ;
\draw [shift={(5.18,4.02)},line width=1.2pt,color=black,fill=blue,fill opacity=0.2]  (0,0) --  plot[domain=2.4965180680191064:3.7478303177768257,variable=\t]({1.*1.8252497086746013*cos(\t r)+0.*1.8252497086746013*sin(\t r)},{0.*1.8252497086746013*cos(\t r)+1.*1.8252497086746013*sin(\t r)}) -- cycle ;
\draw [shift={(5.18,4.02)},line width=1.2pt,color=black,fill=blue,fill opacity=0.2]  (0,0) --  plot[domain=3.7478303177768257:5.002790922933131,variable=\t]({1.*1.8252671037412573*cos(\t r)+0.*1.8252671037412573*sin(\t r)},{0.*1.8252671037412573*cos(\t r)+1.*1.8252671037412573*sin(\t r)}) -- cycle ;
\draw [shift={(5.18,4.02)},line width=1.2pt,color=black,fill=blue,fill opacity=0.2]  (0,0) --  plot[domain=5.002790922933131:6.272196738496853,variable=\t]({1.*1.8160396471443017*cos(\t r)+0.*1.8160396471443017*sin(\t r)},{0.*1.8160396471443017*cos(\t r)+1.*1.8160396471443017*sin(\t r)}) -- cycle ;
\draw [shift={(5.18,4.02)},line width=1.2pt,color=black,fill=blue,fill opacity=0.2]  (0,0) --  plot[domain=-0.010988568682733124:1.241539285032363,variable=\t]({1.*1.8201098867925534*cos(\t r)+0.*1.8201098867925534*sin(\t r)},{0.*1.8201098867925534*cos(\t r)+1.*1.8201098867925534*sin(\t r)}) -- cycle ;
\end{tikzpicture}\hspace{0.5cm}
\begin{tikzpicture}[line cap=round,line join=round,>=triangle 45,x=1.0cm,y=1.0cm,scale=0.7]
\clip(3.,2.) rectangle (7.5,6.);
\draw [line width=1.6pt] (5.18,4.02) circle (1.8201098867925534cm);
\draw [line width=1.2pt] (7.,4.)-- (5.18,4.02);
\draw [line width=1.2pt] (5.18,4.02)-- (5.767189238011339,5.738460372547006);
\draw [line width=1.2pt] (3.7215252327692143,5.117446058978674)-- (5.18,4.02);
\draw [line width=1.2pt] (5.18,4.02)-- (3.68,2.98);
\draw [line width=1.2pt] (5.18,4.02)-- (5.7,2.28);
\draw [shift={(5.18,4.02)},line width=1.2pt,color=black,fill=blue,fill opacity=0.2]  (0,0) --  plot[domain=1.241539285032363:2.4965180680191064,variable=\t]({1.*1.8160113582383595*cos(\t r)+0.*1.8160113582383595*sin(\t r)},{0.*1.8160113582383595*cos(\t r)+1.*1.8160113582383595*sin(\t r)}) -- cycle ;
\draw [shift={(5.18,4.02)},line width=1.2pt,color=black,fill=blue,fill opacity=0.2]  (0,0) --  plot[domain=2.4965180680191064:3.7478303177768257,variable=\t]({1.*1.8252497086746013*cos(\t r)+0.*1.8252497086746013*sin(\t r)},{0.*1.8252497086746013*cos(\t r)+1.*1.8252497086746013*sin(\t r)}) -- cycle ;
\draw [shift={(5.18,4.02)},line width=1.2pt,color=black,fill=blue,fill opacity=0.0]  (0,0) --  plot[domain=3.7478303177768257:5.002790922933131,variable=\t]({1.*1.8252671037412573*cos(\t r)+0.*1.8252671037412573*sin(\t r)},{0.*1.8252671037412573*cos(\t r)+1.*1.8252671037412573*sin(\t r)}) -- cycle ;
\draw [shift={(5.18,4.02)},line width=1.2pt,color=black,fill=blue,fill opacity=0.0]  (0,0) --  plot[domain=5.002790922933131:6.272196738496853,variable=\t]({1.*1.8160396471443017*cos(\t r)+0.*1.8160396471443017*sin(\t r)},{0.*1.8160396471443017*cos(\t r)+1.*1.8160396471443017*sin(\t r)}) -- cycle ;
\draw [shift={(5.18,4.02)},line width=1.2pt,color=black,fill=blue,fill opacity=0.0]  (0,0) --  plot[domain=-0.010988568682733124:1.241539285032363,variable=\t]({1.*1.8201098867925534*cos(\t r)+0.*1.8201098867925534*sin(\t r)},{0.*1.8201098867925534*cos(\t r)+1.*1.8201098867925534*sin(\t r)}) -- cycle ;
\end{tikzpicture}
 
\vspace{-2cm}
\begin{flushright}
\begin{tikzpicture}[x=3.5cm]  
\draw[thick,->] (0,0) -- (2.4,0);
\foreach \x in {0,1,2,}
\draw[shift={(\x,0)},color=black] (0pt,6pt) -- (0pt,-6pt) node [below] {\x};
\foreach \x in {0.2,0.4,0.6,0.8,1.2,1.4,1.6,1.8,2.2}	
\draw[shift={(\x,0)},color=black] (0pt,3pt) -- (0pt,-3pt) ;
\begin{scriptsize}
\draw [color=blue] (1.4,0)-- ++(-3.5pt,0 pt) -- ++(7.0pt,0 pt) ++(-3.5pt,-3.5pt) -- ++(0 pt,7.0pt);
\draw[color=blue] (1.4,-0.8) node {\large{$\dfrac{7}{5}$}};
\end{scriptsize}
\end{tikzpicture}
\end{flushright}


\begin{rmq}
Lorsque le dénominateur est égal à $10$, $100$, $1000$, \dots on dit que c'est une fraction décimale, par exemple $\dfrac{93}{100}$ ou $\dfrac{6}{10}$
\end{rmq}

\begin{meth}[Prendre une fraction d'un quantité]
Pour prendre une fraction d'une quantité, on multiplie la fraction par cette quantité.

\medskip
\noindent Calculer les $\dfrac{2}{5}$ de $400$, c'est calculer :
\begin{enumerate}
\item[•]$\dfrac{2}{5} \times 400 = \dfrac{400}{5} \times 2 = 80 \times 2 = 160$
\item[•]$\dfrac{2}{5} \times 400 = \dfrac{2 \times 400}{5} = \dfrac{800}{5} = 160$
\item[•]$\dfrac{2}{5} \times 400 = 0,4 \times 400 = 160$
\end{enumerate}
\end{meth}

\vspace{-2mm}
\section{Quotients égaux}

\begin{meth}[critère de divisibilité]
\begin{enumerate}
\item[•] Un nombre est divisible par $2$ si son chiffre des unités est pair (il se termine par $0$, $2$, $4$, $6$ ou $8$)
\item[•] Un nombre est divisible par $5$ si son chiffre des unités est $0$ ou $5$.
\item[•] Un nombre est divisible par $3$ si la somme de ses chiffres est divisible par $3$.

Exemple : $249$ est divisible par $3$ car $2 + 4 + 9 = 15$ qui est divisible par $3$ ($3 \times 5 = 15$).

\item[•] Un nombre est divisible par $9$ si la somme de ses chiffres est divisible par $9$.

Exemple : $2 538$ est divisible par $9$ car $2 + 5 + 3 + 8 = 18$ qui est divisible par $9$ ($3 \times 9 = 18$).
\end{enumerate}
\end{meth}

\begin{prop}
Un quotient ne change pas si l'on multiplie (ou divise) son numérateur ET son dénominateur par un même nombre non nul.
\end{prop}

\begin{ex}
\noindent\grandscarreaux{\linewidth}{4cm}
\end{ex}

\subsection{Simplifier une fraction}

\begin{dft}
Simplifier une fraction, c'est écrire une fraction qui lui est égale, mais avec un numérateur et un dénominateur plus petit.
\end{dft}

\begin{ex}
\noindent\grandscarreaux{\linewidth}{3.2cm}
\end{ex}

\begin{rmq}
Aucun autre nombre que $1$ ne divise à la fois $3$ et $4$, la fraction $\dfrac{3}{4}$ ne peut plus être simplifiée. On dit que cette fraction est irréductible.
\end{rmq}

\subsection{Division par un nombre décimal}

\begin{prop}
Pour diviser par un nombre décimal non entier, on se ramène à la division par un nombre entier en multipliant le dividende et le diviseur par $10$ ou par $100$ ou par $1000$ \dots
\end{prop}

\begin{ex}
Calculer $3,5 \div 1,4$ ?

\medskip
\noindent\grandscarreaux{\linewidth}{3.2cm}
\end{ex}

\end{document}